\section{Les règles de sécurité}

\begin{frame}{Les règles de sécurité}
  \begin{block}{Principe général}
    \begin{itemize}
      \item La plongée est une activité de loisir, qui se pratique
            dans le respect strict des règles de sécurité.
      \item Ces règles s’appliquent \textbf{avant, pendant et après la plongée}.
    \end{itemize}
  \end{block}
\end{frame}

\begin{frame}{Avant la plongée}
\begin{itemize}
  \item \alert{Ne jamais plonger si on n’en a pas envie}
        (fatigue, traitements médicaux, affections ORL, gueule de bois…).
  \item Contrôler son matériel :
        \begin{itemize}
          \item pression de la bouteille,
          \item fonctionnement du détendeur,
          \item présence de la ceinture de lest,
          \item gilet et direct system.
        \end{itemize}
  \item Écouter les consignes du \textbf{directeur de plongée}
        et de son \textbf{guide de palanquée}.
  \item \alert{Ne jamais plonger seul}.
\end{itemize}
\end{frame}

\begin{frame}{Mise à l’eau et descente}
\begin{block}{Mise à l’eau}
\begin{itemize}
  \item Le guide de palanquée se met à l’eau en premier.
  \item Gonfler le gilet en surface.
  \item Rester groupé avec tous les membres de sa palanquée.
\end{itemize}
\end{block}

\vspace{0.2cm}

\begin{block}{À la descente}
\begin{itemize}
  \item Descendre uniquement sur ordre du guide de palanquée.
  \item \alert{Ne jamais descendre plus bas que son guide de palanquée}.
  \item Stopper la descente en cas de douleur
        (sinus, oreilles, estomac).
\end{itemize}
\end{block}
\end{frame}

\begin{frame}{Pendant la plongée}
\begin{itemize}
  \item Ne jamais se laisser tenter par la profondeur.
  \item Si un plongeur fait le signe \textbf{« réserve »},
        \alert{toute la palanquée remonte}.
  \item Signaler au guide de palanquée que vous êtes en réserve.
  \item Signaler tout évènement au guide de palanquée.
  \item En cas de perte de la palanquée :
        \begin{itemize}
          \item attendre \textbf{1 minute},
          \item puis remonter à \textbf{10 m/min},
          \item et retour en surface.
        \end{itemize}
\end{itemize}
\end{frame}

\begin{frame}{Remontée et surface}

\begin{block}{À la remontée}
\begin{itemize}
  \item \alert{Ne jamais bloquer sa respiration}.
  \item Ne jamais remonter plus vite que son guide de palanquée.
  \item Respecter la vitesse de remontée de \textbf{10 m/min}
        et les paliers.
  \item Regarder (360°) et écouter avant de faire surface.
\end{itemize}
\end{block}

\vspace{0.2cm}

\begin{block}{En surface}
\begin{itemize}
  \item Attendre toute la palanquée pour rejoindre le bateau.
  \item Attendre que l’échelle soit dégagée.
  \item \alert{Ne pas rester sous l’échelle}.
\end{itemize}
\end{block}
\end{frame}

\begin{frame}{Règles de sécurité — Après la plongée}

\begin{itemize}
  \item Écouter le briefing du guide de palanquée.
  \item \alert{Signaler toute sensation inhabituelle} :
        fatigue, nausée, picotements, malaise…
  \item Ne pas prendre l’avion dans les \textbf{12 à 24 heures}
        suivant la plongée.
  \item Ne pas faire d'effort physique intense
        dans les \textbf{6 heures} suivant la plongée.
  \item Ne pas faire de l'apnée dans les \textbf{6 heures} suivant la plongée.
\end{itemize}
\end{frame}
